\chapter{Schlussfolgerung}
\label{chap:Schlussfolgerung}
Ziel meiner Arbeit war es, eine funktionsfähige Messenger-Software zu entwickeln, die verschlüsselte Nachrichten über ein selbst aufgebautes Protokoll verschickt. Bei der Umsetzung stiess ich auf mehrere Einschränkungen.
Ein vollständig eigenes Protokoll liess sich nicht realisieren. Alle Geräte kommunizieren über standardisierte Protokolle, und die Implementierung eines eigenen Protokolls hätte mehr Zeit beansprucht, als im Rahmen einer Maturaarbeit zur Verfügung steht. Zudem könnten stark veränderte Datenpakete vom Schulnetzwerk aus Sicherheitsgründen automatisch blockiert werden. Mein Protokoll betrifft deshalb nur die Nutzdaten (Payload) und wird vom TCP/IP-Protokoll gekapselt, das die Teilnehmer miteinander verbindet.
Die Implementierung entspricht ausserdem nicht dem aktuellen Stand der Technik. Sie kann Sicherheitslücken enthalten, die mir nicht bekannt sind und die sich leicht ausnutzen liessen. Damit könnten alle Benutzer gefährdet sein.
Funktional ist die Software auf Textnachrichten beschränkt, der Versand von Bildern oder Dateien wird nicht unterstützt. Um eine Verbindung aufzubauen, müssen die Verbindungsparameter in den Programmdateien manuell angepasst werden. Das Programm läuft mit Python und dem Werkzeug uv und wird über die Kommandozeile gestartet. Die grafische Benutzeroberfläche ist dagegen übersichtlich aufgebaut und einfach zu bedienen.
Ein weiterer Kritikpunkt ist, dass das Projekt keine Threading Locks verwendet, obwohl sie in der Theorie behandelt werden. Dadurch können beim Lesen und Überschreiben der Client-Liste Fehler entstehen, wenn zwei Benutzer fast gleichzeitig Änderungen vornehmen.
Ausserdem geht die Implementierung davon aus, dass ein Aufruf von \texttt{recv()} immer eine vollständige Nachricht liefert. TCP ist jedoch ein Datenstrom ohne feste Nachrichtengrenzen, deshalb könnten Nachrichten bei hoher Netzwerkauslastung unvollständig oder zusammengefasst ankommen.
Die AES-Verschlüsselung im EAX-Modus schützt zwar den Inhalt der Nachrichten vor Angreifern, gibt aber IP-Adresse und Port des Empfängers im Klartext preis. Ein Angreifer kann dadurch nachvollziehen, wer mit wem kommuniziert.